Let $\mathcal{P}$ be a projective plane of order $n$ and let $R$ be any set of symbols with cardinal $n$ such that $0, 1 \in R$, $0 \neq 1$, but the symbol ``$\infty$'' is not in $R$. Choose any line of $\mathcal{P}$ and label this line $\ell_0$ and then pick any other two lines $\ell_1, \ell_2$ with the only restriction that $\ell_1, \ell_2, \ell_0$ form the sides of a non-degenerate triangle. Let $X = \ell_2 \ell_0$, $Y = \ell_0 \ell_1$ and $O = \ell_1 \ell_2$. Finally let $I$ be any point not incident with any side of the chosen triangle.

\textbf{Coordinatization of points.} Use the elements of $R$ plus the extra symbol $\infty$ to coordinatize $\mathcal{P}$ relative to the quadrangle $O, X, Y, I$. Label $A = XI \cap \ell_1$, $B = YI \cap \ell_2$, and $J = AB \cap XY$ (where $XY = \ell_0$). Assign coordinates to the points of $\mathcal{P}$ as follows:

\begin{itemize}
\item Label $Y$ as $(\infty)$.
\item Label the point $OI \cap XY$ as $(1)$; all other points on $\ell_0$ except $Y$ receive distinct symbols from $R$ as coordinates $(m)$.
\item On $\ell_1$, label $O$ as $(0,0)$, $A$ as $(0,1)$, $Y$ as $(0)$ (which is the same as $(\infty)$), and distribute the remaining symbols of $R$ as $(0,a)$ to the other points of $\ell_1 \setminus Y$.
\item For any point $D$ on $\ell_2$, $D \neq X$, let $D' = JD \cap \ell_1$ and, if $D'$ is $(0,d)$, write $(d,0)$ for $D$. Thus $O$ has coordinates $(0,0)$.
\item For any point $E$ not on $\ell_\infty$ (which is $\ell_0$), if $XE \cap \ell_1$ is $(0,g)$ and $YE \cap \ell_2$ is $(f,0)$, then $E$ is given the coordinates $(f,g)$.
\item If $M$ is any point of $\ell_\infty$ other than $Y$ and if the line joining $M$ to $(1,0)$ meets $\ell_1$ in $(0,m)$, then $M$ is given the coordinate $(m)$.
\end{itemize}

\textbf{Coordinatization of lines.} If $\ell$ is any line not containing $Y$, then, if $\ell$ meets $\ell_\infty$ at the point $(m)$ and $\ell_1$ at the point $(0,k)$, give $\ell$ the coordinates $[m,k]$. If $\ell$ contains $Y$ but is distinct from $\ell_\infty$, then call $\ell$ the line $[k]$ where $k$ is determined by $\ell \cap \ell_2 = (k,0)$. Finally call $\ell_\infty$ the line $[\infty]$.

In this way every point and line of $\mathcal{P}$ has been coordinatized, the coordinatization depending only on the initial choice of $O, X, Y, I$ and the way in which the elements of $R$ are assigned to the points of $\ell_1 \setminus Y$.
